\section{Peer-to-Peer Network Protocol}
\label{sec:p2p}

The P2P protocol transports blocks, transactions, headers, addresses, filters,
and relay-control messages between nodes. This section specifies interoperable
wire behavior.

\subsection{Message envelope}

Legacy unencrypted Bitcoin P2P messages use a 24-byte header followed by a
payload:

\begin{longtable}{@{}L{0.22\linewidth}L{0.20\linewidth}L{0.46\linewidth}@{}}
\toprule
Field & Type & Meaning \\
\midrule
\code{start} & \code{byte[4]} & Network magic identifying mainnet, testnet, signet, or regtest. \\
\code{command} & \code{byte[12]} & ASCII command name padded with zero bytes. \\
\code{payloadSize} & \code{uint32} & Number of payload bytes. \\
\code{checksum} & \code{byte[4]} & First four bytes of \(\mathrm{H}(payload)\). \\
\code{payload} & \code{byte[]} & Command-specific serialized payload of length \code{payloadSize}. \\
\bottomrule
\end{longtable}

BIP324 defines an encrypted v2 transport that changes transport framing after
negotiation, but command payload semantics remain command-specific
\cite{bip-0324,bitcoin-core-v31}.

\subsection{Handshake}

A connection becomes usable after version negotiation:

\begin{enumerate}[leftmargin=*]
  \item each side sends \code{version};
  \item each side validates network, protocol version, services, timestamp,
        nonce, and user-agent limits;
  \item each side sends \code{verack}; and
  \item optional relay, compact-block, address-relay, filter, and transport
        features are negotiated by later messages.
\end{enumerate}

\begin{longtable}{@{}L{0.24\linewidth}L{0.24\linewidth}L{0.40\linewidth}@{}}
\toprule
\code{version} field & Type & Meaning \\
\midrule
\code{version} & \code{int32} & Protocol version offered by the sender. \\
\code{services} & \code{uint64} & Service bits offered by the sender. \\
\code{timestamp} & \code{int64} & Sender Unix time in seconds. \\
\code{addrRecv} & network address & Address of the receiving peer as understood by the sender. \\
\code{addrFrom} & network address & Address of the sender. \\
\code{nonce} & \code{uint64} & Random connection nonce used to detect self-connections. \\
\code{userAgent} & \code{string} & BIP14 user-agent string. \\
\code{startHeight} & \code{int32} & Sender's current best height estimate. \\
\code{relay} & \code{bool}, optional & BIP37 transaction-inventory relay preference when present. \\
\bottomrule
\end{longtable}

\subsection{Control and negotiation payloads}

\begin{longtable}{@{}L{0.24\linewidth}L{0.24\linewidth}L{0.40\linewidth}@{}}
\toprule
Message & Payload & Meaning \\
\midrule
\code{verack} & empty & Completes version negotiation. \\
\code{sendheaders} & empty & Requests direct \code{headers} announcements instead of block inventory announcements. \\
\code{wtxidrelay} & empty & Negotiates transaction relay by witness transaction identifier. \\
\code{sendaddrv2} & empty & Negotiates BIP155 address relay. \\
\code{ping} & \code{uint64 nonce} & Liveness probe. \\
\code{pong} & \code{uint64 nonce} & Response echoing the received ping nonce. \\
\code{feefilter} & \code{uint64 feerate} & Asks the peer not to announce transactions below the feerate in satoshis per 1000 virtual bytes. \\
\bottomrule
\end{longtable}

\subsection{Service flags}

\begin{longtable}{@{}L{0.24\linewidth}L{0.14\linewidth}L{0.50\linewidth}@{}}
\toprule
Flag & Bit & Meaning \\
\midrule
\code{NODE\_NETWORK} & 0 & Peer can serve the full block chain. \\
\code{NODE\_BLOOM} & 2 & Peer supports BIP37 bloom-filter service. \\
\code{NODE\_WITNESS} & 3 & Peer supports witness transaction and block relay. \\
\code{NODE\_COMPACT\_FILTERS} & 6 & Peer supports BIP157 compact block filter service. \\
\code{NODE\_NETWORK\_LIMITED} & 10 & Peer can serve at least a limited recent block window. \\
\bottomrule
\end{longtable}

Service flags advertise capability. A node must still validate responses and
apply local resource limits.

\subsection{Inventory and object retrieval}

Inventory relay announces objects by type and hash. A peer requests announced
objects with \code{getdata}; the responder sends the corresponding object or a
\code{notfound} response when unavailable.

\begin{longtable}{@{}L{0.20\linewidth}L{0.24\linewidth}L{0.44\linewidth}@{}}
\toprule
Record & Fields & Meaning \\
\midrule
\code{InventoryVector} & \code{type:uint32}, \code{hash:uint256} & Identifies a transaction, block, filtered block, compact block, witness transaction, or witness block. \\
\code{inv} & {\small\code{vector<}\newline\code{InventoryVector>}} & Announces available objects. \\
\code{getdata} & {\small\code{vector<}\newline\code{InventoryVector>}} & Requests announced objects. \\
\code{notfound} & {\small\code{vector<}\newline\code{InventoryVector>}} & Reports requested objects not available from the responder. \\
\bottomrule
\end{longtable}

Witness inventory types request serialization that includes witness data. A
peer that has not negotiated witness support must not rely on receiving witness
objects from that connection \cite{bip-0144}.

\begin{longtable}{@{}L{0.30\linewidth}L{0.20\linewidth}L{0.38\linewidth}@{}}
\toprule
Inventory type & Value & Object requested or announced \\
\midrule
\code{MSG\_TX} & 1 & Transaction by \code{txid}. \\
\code{MSG\_BLOCK} & 2 & Block without requiring witness serialization. \\
\code{MSG\_FILTERED\_BLOCK} & 3 & BIP37 filtered block. \\
\code{MSG\_CMPCT\_BLOCK} & 4 & BIP152 compact block. \\
\code{MSG\_WITNESS\_TX} & $2^{30} + 1$ & Transaction including witness data, identified by \code{txid} request type with witness flag. \\
\code{MSG\_WITNESS\_BLOCK} & $2^{30} + 2$ & Block including witness data. \\
{\small\code{MSG\_FILTERED}\newline\code{\_WITNESS\_BLOCK}} & $2^{30} + 3$ & Filtered block including witness data. \\
\bottomrule
\end{longtable}

\subsection{Header and block synchronization}

Headers synchronize chain work before full block data. A locator is a sequence
of block hashes, newest first, stepping back from a known tip toward genesis.

Messages \code{getheaders} and \code{getblocks} carry a locator and stop hash.
Message \code{headers} returns block headers in forward order, each followed by
CompactSize zero. Message \code{block} carries serialized full block data. A
node must validate proof of work, linkage, target, timestamp, and deployment
context for received headers before using them for chain selection.

\begin{longtable}{@{}L{0.24\linewidth}L{0.24\linewidth}L{0.40\linewidth}@{}}
\toprule
Record & Fields & Encoding rule \\
\midrule
\code{BlockLocator} & \code{version}, \code{hashes} & \code{int32} version followed by \code{vector<uint256>} locator hashes. \\
{\small\code{getheaders}/\newline\code{getblocks}} & \code{version}, \code{locator}, \code{stopHash} & \code{int32}, locator vector, then \code{uint256} stop hash. \\
\code{headers} item & \code{header}, \code{txCount} & 80-byte block header followed by CompactSize zero. A nonzero count is invalid for \code{headers}. \\
\bottomrule
\end{longtable}

\subsection{Compact block relay}

BIP152 compact blocks reduce block-relay bandwidth by sending a header, short
transaction identifiers, selected prefilled transactions, and requesting any
missing transactions \cite{bip-0152}.

Message \code{sendcmpct} negotiates compact-block mode and announcement
preference. Messages \code{cmpctblock}, \code{getblocktxn}, and \code{blocktxn}
carry the compact block and missing-transaction exchange. Compact
reconstruction is an optimization only. The reconstructed full block must pass
the same block validation rules as any received \code{block} message.

\begin{longtable}{@{}L{0.24\linewidth}L{0.24\linewidth}L{0.40\linewidth}@{}}
\toprule
Record & Fields & Encoding rule \\
\midrule
{\small\code{Compact}\newline\code{Block}} & \code{header}, \code{nonce}, \code{shortIds}, \code{prefilledTx} & Header, \code{uint64} nonce, CompactSize count plus 6-byte short ids, then \code{vector<PrefilledTx>}. \\
\code{PrefilledTx} & \code{indexDelta}, \code{tx} & CompactSize differential index from the previous prefilled transaction, then full transaction. \\
\code{getblocktxn} & \code{blockHash}, \code{indexes} & Block hash followed by CompactSize differential transaction indexes. \\
\code{blocktxn} & \code{blockHash}, \code{transactions} & Block hash followed by transaction vector. \\
\bottomrule
\end{longtable}

\subsection{Address relay}

Legacy address records and BIP155 address records use different encodings:

\begin{longtable}{@{}L{0.22\linewidth}L{0.24\linewidth}L{0.42\linewidth}@{}}
\toprule
Record & Fields & Encoding rule \\
\midrule
{\small\code{NetAddress}\newline\code{NoTime}} & \code{services}, \code{address}, \code{port} & \code{uint64} services, 16 address bytes, then \code{uint16} port in big-endian network byte order. Used inside \code{version}. \\
{\small\code{NetAddress}\newline\code{Legacy}} & \code{time}, \code{services}, \code{address}, \code{port} & \code{uint32} time plus \code{NetAddressNoTime}. Used by \code{addr}. \\
\code{AddrV2} & \code{time}, \code{services}, \code{networkId}, \code{address}, \code{port} & \code{uint32} time, CompactSize services, \code{uint8} network id, CompactSize address length, address bytes, then big-endian port. \\
\bottomrule
\end{longtable}

Messages \code{addr} and \code{addrv2} relay legacy and BIP155 address
records. Message \code{sendaddrv2} announces BIP155 address-relay support, and
\code{getaddr} requests known peer addresses subject to local rate limits.
Address relay is unauthenticated gossip. Received addresses are hints, not
validated reachability claims \cite{bip-0155}.

\subsection{Light-client services}

Bitcoin P2P includes two optional light-client service families:

\begin{longtable}{@{}L{0.24\linewidth}L{0.26\linewidth}L{0.38\linewidth}@{}}
\toprule
Service & Messages & Status \\
\midrule
BIP37 bloom filters & \code{filterload}, \code{filteradd}, \code{filterclear}, \code{merkleblock} & Optional peer service; disabled unless a node chooses to provide bloom-filter support. \\
BIP157/158 compact filters & \code{getcfilters}, \code{cfilter}, \code{getcfheaders}, \code{cfheaders}, \code{getcfcheckpt}, \code{cfcheckpt} & Optional compact-filter service for clients that verify filters against block-filter headers. \\
\bottomrule
\end{longtable}

These services do not change block or transaction validity
\cite{bip-0037,bip-0157,bip-0158}.

\begin{longtable}{@{}L{0.24\linewidth}L{0.28\linewidth}L{0.36\linewidth}@{}}
\toprule
Message & Payload fields & Encoding rule \\
\midrule
\code{filterload} & \code{filter}, \code{hashFuncs}, \code{tweak}, \code{flags} & \code{varbytes}, \code{uint32}, \code{uint32}, \code{uint8}. \\
\code{filteradd} & \code{data} & \code{varbytes}. \\
\code{filterclear} & empty & Clears the peer's BIP37 filter state. \\
\code{merkleblock} & \code{header}, \code{total}, \code{hashes}, \code{flags} & Block header, \code{uint32} total transactions, \code{vector<uint256>} hashes, and flag bytes. \\
\code{getcfilters} & \code{type}, \code{startHeight}, \code{stopHash} & \code{uint8}, \code{uint32}, \code{uint256}. \\
\code{cfilter} & \code{type}, \code{blockHash}, \code{filter} & \code{uint8}, \code{uint256}, \code{varbytes}. \\
\code{getcfheaders} & \code{type}, \code{startHeight}, \code{stopHash} & \code{uint8}, \code{uint32}, \code{uint256}. \\
\code{cfheaders} & \code{type}, \code{stopHash}, \code{prevHeader}, \code{filterHashes} & \code{uint8}, two \code{uint256} values, then \code{vector<uint256>}. \\
\code{getcfcheckpt} & \code{type}, \code{stopHash} & \code{uint8}, \code{uint256}. \\
\code{cfcheckpt} & \code{type}, \code{stopHash}, \code{headers} & \code{uint8}, \code{uint256}, \code{vector<uint256>}. \\
\bottomrule
\end{longtable}
